% !TeX root = ../main.tex
% 致谢中主要感谢导师和对论文工作有直接贡献和帮助的人士和单位。

% 一般致谢的内容有：
% \begin{enumerate}[label=（\chinese*）,itemindent=2em]

%     \item 对指导或协助指导完成论文的导师；
%     \item 对国家科学基金、资助研究工作的奖学金基金、合同单位、资助或支持的企业、组织或个人；
%     \item 对协助完成研究工作和提供便利条件的组织或个人；
%     \item 对在研究工作中提出建议和提供帮助的人；
%     \item 对给予转载和引用权的资料、图片、文献、研究思想和设想的所有者；
%     \item 对其他应感谢的组织和个人。

% \end{enumerate}

% 致谢言语应谦虚诚恳，实事求是。字数不超过1000汉字

随着这篇论文的最终定稿，我的硕士研究生生涯乃至整个漫长的学生时代，也即将画上圆满的句号。回首在校园里度过的这无数个日日夜夜，心中感慨万千，既有对知识殿堂的不舍，也有对未来旅途的期盼。在此，我想向所有在此期间给予我帮助与支持的人，致以最诚挚的感谢。

首先，我要向我的导师们表达最深切的敬意与谢意。感谢实验室带头人刘均老师，您高屋建瓴的学术视野总是在大方向上为我指引道路，让我得以在宏观上把握研究的脉络。特别要感谢我的导师张玲玲老师，从课题的选定、具体方向的深入探索，到这篇论文的字斟句酌与反复修改，都得益于张老师背后的指导和鼓励。同时，我也要感谢魏笔凡老师和马杰老师，感谢你们在历次组会上给予我的专业指导与中肯建议，让我受益良多。在此，祝愿各位老师身体健康，工作顺利，家庭幸福！

其次，我要感谢在实验室并肩作战的同学们。感谢同级同学：茹燕池、郭天霖、张宇舜、孙辰扬、陈一心、朱隆吉、唐熙、渠宁、赵宏涛、黄家兴、王庆康，我们一起探讨问题、互相打气的美好时光，我将永远铭记心中。也要特别感谢付雨萌师妹对本论文的支持。最后，我也要感谢其他所有为本论文提出过宝贵意见的老师和同学，我也很感谢看到这段致谢的你。

此外，我要深深感谢我的家人。感谢你们二十多年来如一日的无私奉献与默默付出。是你们在背后的包容、理解与支持，给了我不断前行的底气和最温暖的避风港，让我能够心无旁骛地完成学业。

韶华易逝，岁月流金。即将挥手告别这座充满青春记忆的校园，去往未知的人生下一站，心中难免有对过往的眷恋，也有对新生活的憧憬与迷茫。未来的路或许平坦，或许泥泞，正如歌词：你走在这繁华的街上，在寻找你该去的方向，你走在这繁华的街上，在寻找你曾拥有的力量。愿大家能带着在校园里积蓄的这份力量，勇敢地走向人生的下一个远方。

% \vspace{\baselineskip}
% {\color{red} 用于双盲评审的论文，此页内容全部隐去。}

